\documentclass[tikz,border=8pt]{standalone}
\usepackage{tikz}
\usepackage{amsmath}
\usepackage{amssymb}
\usepackage{pgfplots}
\pgfplotsset{compat=1.18}
\usepackage{ctex}
\usetikzlibrary{calc,positioning,arrows.meta,matrix}

% 最大堆 vs 最小堆结构对比

\begin{document}
\begin{tikzpicture}[
  every node/.style={font=\small},
  maxnode/.style={draw=orange!70, circle,
    minimum size=0.8cm, thick, fill=orange!10},
  minnode/.style={draw=blue!60, circle,
    minimum size=0.8cm, thick, fill=blue!8},
  cell/.style={draw=gray!60, rectangle,
    minimum width=0.75cm, minimum height=0.65cm, thick, fill=white},
  arr/.style={->,>=Stealth,gray!60,thick},
]

%% ── 标题 ──────────────────────────────────────────────────
\node[font=\large\bfseries] at (9.0, 1.2)
  {堆结构：最大堆 vs 最小堆};
\node[font=\small,gray] at (9.0, 0.55)
  {完全二叉树存储于数组；父节点下标 \texttt{(i-1)//2}，子节点 \texttt{2i+1}, \texttt{2i+2}};

%% ══════════════════════════════════════════════════════════
%% 左侧：最大堆
%% 树结构：
%%       10 (idx 0)
%%      /  \
%%     7    9  (idx 1,2)
%%    / \  / \
%%   4  6 3  8  (idx 3,4,5,6)
%% ══════════════════════════════════════════════════════════

\def\lx{3.5}  % 最大堆中心 x

\node[font=\small\bfseries,orange!70] at (\lx, -0.2) {最大堆（Max-Heap）};
\node[font=\tiny,orange!60] at (\lx, -0.65) {父节点 $\geq$ 子节点};

% 层 0
\node[maxnode] (mx0) at (\lx, -1.5) {10};

% 层 1
\node[maxnode] (mx1) at (\lx-1.4, -3.0) {7};
\node[maxnode] (mx2) at (\lx+1.4, -3.0) {9};

% 层 2
\node[maxnode] (mx3) at (\lx-2.1, -4.5) {4};
\node[maxnode] (mx4) at (\lx-0.7, -4.5) {6};
\node[maxnode] (mx5) at (\lx+0.7, -4.5) {3};
\node[maxnode] (mx6) at (\lx+2.1, -4.5) {8};

% 边
\draw[arr] (mx0) -- (mx1);
\draw[arr] (mx0) -- (mx2);
\draw[arr] (mx1) -- (mx3);
\draw[arr] (mx1) -- (mx4);
\draw[arr] (mx2) -- (mx5);
\draw[arr] (mx2) -- (mx6);

% 父≥子标注（根旁）
\node[font=\tiny,orange!70,anchor=west] at (\lx+0.45, -1.5) {根=最大};

% 数组表示
\node[font=\scriptsize,gray] at (\lx, -5.3) {数组表示：};
\foreach \i/\v in {0/10, 1/7, 2/9, 3/4, 4/6, 5/3, 6/8} {
  \node[cell,font=\small] at (\lx-2.25+\i*0.75, -5.9) {\v};
  \node[font=\tiny,gray] at (\lx-2.25+\i*0.75, -6.4) {\i};
}

% 节点与数组对应连线（只标根和末尾）
\draw[->,>=Stealth,orange!30,very thin,dashed]
  (mx0.south) -- ++(0,-0.5) -| (\lx-2.25, -5.65);

%% ── 父子关系说明（最大堆）────────────────────────────────
\node[draw=gray!40,fill=white,thin,rounded corners=2pt,
      font=\scriptsize,inner sep=4pt,align=left]
  at (\lx, -7.2)
  {idx 0: 10（根，最大）\\
   idx 1: 7，idx 2: 9（10的子）\\
   idx 3: 4，idx 4: 6（7的子）\\
   idx 5: 3，idx 6: 8（9的子）\\
   性质：\texttt{heap[i] $\geq$ heap[2i+1], heap[2i+2]}};

%% ══════════════════════════════════════════════════════════
%% 右侧：最小堆
%% 树结构：
%%        1 (idx 0)
%%       / \
%%      3   2 (idx 1,2)
%%     / \ / \
%%    7  4 9  5 (idx 3,4,5,6)
%% ══════════════════════════════════════════════════════════

\def\rx{14.0}  % 最小堆中心 x

\node[font=\small\bfseries,blue!70] at (\rx, -0.2) {最小堆（Min-Heap）};
\node[font=\tiny,blue!60] at (\rx, -0.65) {父节点 $\leq$ 子节点};

% 层 0
\node[minnode] (mn0) at (\rx, -1.5) {1};

% 层 1
\node[minnode] (mn1) at (\rx-1.4, -3.0) {3};
\node[minnode] (mn2) at (\rx+1.4, -3.0) {2};

% 层 2
\node[minnode] (mn3) at (\rx-2.1, -4.5) {7};
\node[minnode] (mn4) at (\rx-0.7, -4.5) {4};
\node[minnode] (mn5) at (\rx+0.7, -4.5) {9};
\node[minnode] (mn6) at (\rx+2.1, -4.5) {5};

% 边
\draw[arr] (mn0) -- (mn1);
\draw[arr] (mn0) -- (mn2);
\draw[arr] (mn1) -- (mn3);
\draw[arr] (mn1) -- (mn4);
\draw[arr] (mn2) -- (mn5);
\draw[arr] (mn2) -- (mn6);

% 标注
\node[font=\tiny,blue!70,anchor=west] at (\rx+0.45, -1.5) {根=最小};

% 数组表示
\node[font=\scriptsize,gray] at (\rx, -5.3) {数组表示：};
\foreach \i/\v in {0/1, 1/3, 2/2, 3/7, 4/4, 5/9, 6/5} {
  \node[cell,font=\small] at (\rx-2.25+\i*0.75, -5.9) {\v};
  \node[font=\tiny,gray] at (\rx-2.25+\i*0.75, -6.4) {\i};
}

\draw[->,>=Stealth,blue!30,very thin,dashed]
  (mn0.south) -- ++(0,-0.5) -| (\rx-2.25, -5.65);

%% ── 父子关系说明（最小堆）────────────────────────────────
\node[draw=gray!40,fill=white,thin,rounded corners=2pt,
      font=\scriptsize,inner sep=4pt,align=left]
  at (\rx, -7.2)
  {idx 0: 1（根，最小）\\
   idx 1: 3，idx 2: 2（1的子）\\
   idx 3: 7，idx 4: 4（3的子）\\
   idx 5: 9，idx 6: 5（2的子）\\
   性质：\texttt{heap[i] $\leq$ heap[2i+1], heap[2i+2]}};

%% ── 中间分隔线 ───────────────────────────────────────────
\draw[gray!30,dashed,thick] (9.0, -0.0) -- (9.0, -8.5);
\node[font=\scriptsize,gray] at (9.0, -8.8) {vs};

%% ── 底部分隔线 ───────────────────────────────────────────
\draw[gray!25,dashed] (0.0, -9.2) -- (18.0, -9.2);

%% ── 算法说明 ─────────────────────────────────────────────
\node[draw=gray!50,fill=white,rounded corners=3pt,font=\small,
      inner sep=8pt,align=left]
  at (9.0, -10.3)
  {\textbf{堆操作复杂度：}\\[4pt]
   \textbf{插入（push）：}$O(\log n)$，加到末尾后上浮（sift-up）\\[2pt]
   \textbf{弹出（pop）：}$O(\log n)$，取出根后将末尾移到根再下沉（sift-down）\\[2pt]
   \textbf{建堆（heapify）：}$O(n)$，从最后一个非叶节点向上执行 sift-down\\[2pt]
   Python: \texttt{heapq}（最小堆）；最大堆取反：\texttt{heappush(h, -x)}};

\end{tikzpicture}
\end{document}
